% ============================================================
% EJERCICIO 2 — Tablas y datos experimentales
% Objetivo : tabla profesional + modelo + cálculo de tau.
% Datos    : t (ms) = 0,10,20,30,40,50 ; v (V) = 0; 2,1; 3,5; 4,3; 4,7; 4,9
% Requisitos: tabla con booktabs, título y etiqueta; ecuación del modelo
%            v_C(t)=V_in(1-e^{-t/tau}) con V_in = 5 V; tau estimado
%            (t en que v_C = 0,632 V_in); gráfica con pgfplots (opcional +2).
% PROMPT UTILIZADO: TODO
% ============================================================
\documentclass[12pt,a4paper]{article}
\usepackage[utf8]{inputenc}
\usepackage[T1]{fontenc}
\usepackage[spanish]{babel}
\usepackage{amsmath,amssymb}
\usepackage{booktabs}
\usepackage{siunitx}
\sisetup{locale=FR, output-decimal-marker={,}}
\usepackage{circuitikz}
\usepackage[hidelinks]{hyperref}

\begin{document}
\section{Datos experimentales}
% TODO tabla: \begin{table}[htbp] ... \toprule ... \caption{...} \label{tab:carga}

\section{Modelo}
% TODO ecuación numerada con \label{eq:carga}

\section{Estimación de $\tau$}
% TODO: 0,632 \cdot 5 V = 3,16 V. Interpole entre dos filas de la Tabla~\ref{tab:carga}.

\end{document}
% RÚBRICA (20 pts): tabla sin líneas verticales y con caption/label (5)
% · ecuación correcta (4) · tau estimado y justificado (5)
% · referencias \ref/\eqref correctas (3) · unidades con siunitx (3)
